%% 
%% Copyright 2007-2020 Elsevier Ltd
%% 
%% This file is part of the 'Elsarticle Bundle'.
%% ---------------------------------------------
%% 
%% It may be distributed under the conditions of the LaTeX Project Public
%% License, either version 1.2 of this license or (at your option) any
%% later version.  The latest version of this license is in
%%    http://www.latex-project.org/lppl.txt
%% and version 1.2 or later is part of all distributions of LaTeX
%% version 1999/12/01 or later.
%% 
%% The list of all files belonging to the 'Elsarticle Bundle' is
%% given in the file `manifest.txt'.
%% 
%% Template article for Elsevier's document class `elsarticle'
%% with harvard style bibliographic references

\documentclass[preprint,12pt,authoryear]{elsarticle}

%% Use the option review to obtain double line spacing
%% \documentclass[authoryear,preprint,review,12pt]{elsarticle}

%% Use the options 1p,twocolumn; 3p; 3p,twocolumn; 5p; or 5p,twocolumn
%% for a journal layout:
%% \documentclass[final,1p,times,authoryear]{elsarticle}
%% \documentclass[final,1p,times,twocolumn,authoryear]{elsarticle}
%% \documentclass[final,3p,times,authoryear]{elsarticle}
%% \documentclass[final,3p,times,twocolumn,authoryear]{elsarticle}
%% \documentclass[final,5p,times,authoryear]{elsarticle}
%% \documentclass[final,5p,times,twocolumn,authoryear]{elsarticle}

%% For including figures, graphicx.sty has been loaded in
%% elsarticle.cls. If you prefer to use the old commands
%% please give \usepackage{epsfig}

%% The amssymb package provides various useful mathematical symbols
\usepackage{amssymb}
\usepackage{cite}
\usepackage{amsmath,amssymb,amsfonts}
\usepackage{algorithmic}
\usepackage{graphicx}
\usepackage{textcomp}
\usepackage{graphics}
\usepackage{caption}

\usepackage{listings}
\usepackage{float}
%%%%%%%%%%%%%%%%%%%
\usepackage{longtable}
\usepackage{booktabs}
\usepackage{lscape} % Para orientación de página horizontal
\usepackage{graphicx} % Para manejar imágenes y gráficos
\usepackage{supertabular}
%% The amsthm package provides extended theorem environments
%% \usepackage{amsthm}

%% The lineno packages adds line numbers. Start line numbering with
%% \begin{linenumbers}, end it with \end{linenumbers}. Or switch it on
%% for the whole article with \linenumbers.
%% \usepackage{lineno}

\journal{Accident Analysis \& Prevention}

\begin{document}

\begin{frontmatter}

%% Title, authors and addresses

%% use the tnoteref command within \title for footnotes;
%% use the tnotetext command for theassociated footnote;
%% use the fnref command within \author or \affiliation for footnotes;
%% use the fntext command for theassociated footnote;
%% use the corref command within \author for corresponding author footnotes;
%% use the cortext command for theassociated footnote;
%% use the ead command for the email address,
%% and the form \ead[url] for the home page:
%% \title{Title\tnoteref{label1}}
%% \tnotetext[label1]{}
\author[1]{Ramon A. Briseño}
\ead{ramon.briseno@academicos.udg.mx}
%% \ead[url]{home page}
%% \fntext[label2]{}
%% \cortext[cor1]{}
\affiliation[1]{organization={Universidad De Guadalajara, CUGDL},
            addressline={Calle Guanajuato. Núm. 1045}, 
            city={Guadalajara},
%%            postcode={}, 
            state={Jalisco},
            country={México}}
%% \fntext[label3]{}
\author[2]{Edgar Cossio}
\ead{edgar.cossio@iieg.gob.mx}
%% \ead[url]{home page}
%% \fntext[label2]{}
%% \cortext[cor1]{}
\affiliation[2]{organization={Instituto de Información Estadística y Geográfica de Jalisco},
            addressline={Calzada de los Pirules 71 Col. Cd. Granja}, 
            city={Zapopan},
            postcode={45010}, 
            state={Jalisco},
            country={México}}
%% \fntext[label3]{}
\author[3]{Carolina Del-Valle-Soto}
\ead{cvalle@up.edu.mx}
%% \ead[url]{home page}
%% \fntext[label2]{}
 \cortext[cor1]{corresponding author}
\affiliation[3]{organization={Universidad Panamericana, Facultad de Ingeniería},
            addressline={Álvaro del Portillo 49}, 
            city={Zapopan},
            postcode={45010}, 
            state={Jalisco},
            country={México}}
%% \fntext[label3]{}
\author[4]{Carlos Mex-Perera}
\ead{jnolazco@tec.mx}
%% \ead[url]{home page}
%% \fntext[label2]{}
%% \cortext[cor1]{}
\affiliation[4]{organization={Rochester Institute of Technology},
            addressline={One Lomb Memorial Dr}, 
            city={Rochester}, 
            state={New York},
            country={USA}}
%% \fntext[label3]{}

\title{Real-time platform for predicting the severity of cyclist-vehicle  accidents in a georeferenced manner using machine learning algorithms}

%% use optional labels to link authors explicitly to addresses:
%% \author[label1,label2]{}
%% \affiliation[label1]{organization={},
%%             addressline={},
%%             city={},
%%             postcode={},
%%             state={},
%%             country={}}
%%
%% \affiliation[label2]{organization={},
%%             addressline={},
%%             city={},
%%             postcode={},
%%             state={},
%%             country={}}


\begin{abstract}
%% Text of abstract
To reduce the mobility carbon footprint in megacities, bicycles for long distances is one of the clean and healthy solution. However, the amount of registered accidents on the road between motorized vehicles and bicycles creates an important concern. Moreover, we present the case of the Guadalajara Metropolitan Area with more than 6 million population, 3 million vehicles, and in the past decade more than 120km of bicycle lanes looking to provide safety to more than 126,000 bicycle users. For that achievement, we can use the historical registries of bicycle accidents to train machine learning algorithms in order to predict the severity of cyclist-vehicle accidents, identify dangerous crossroads under certain conditions, and inform users into interactive mobility maps to reduce possibilities of accidents. We call that set of machine learning algorithms the RAMCI platform, and this work presents how this approach is more efficient than other similar systems developed in the world to predict the severity of cyclist-vehicle accidents. The main contribution of this work is the methodology in order to train with the city data, the machine learning algorithms to predict the severity of cyclist-vehicle accidents concerns in city maps with a compressible trusty level. The results achieved, can help other cities looking to develop bicycle mobility as a strategy to reduce carbon footprint to reduce possibility of accidents and to improve adoption of this transport system.
\end{abstract}

%%Graphical abstract
%\begin{graphicalabstract}
%\includegraphics{grabs}
%\end{graphicalabstract}

%%Research highlights
%\begin{highlights}
%\item Research highlight 1
%\item Research highlight 2
%\end{highlights}

\begin{keyword}
%% keywords here, in the form: keyword \sep keyword
Cyclist-vehicle accident \sep bicycle mobility \sep real-time platform \sep severity prediction \sep machine learning.
%% PACS codes here, in the form: \PACS code \sep code

%% MSC codes here, in the form: \MSC code \sep code
%% or \MSC[2008] code \sep code (2000 is the default)

\end{keyword}

\end{frontmatter}

%% \linenumbers

%% main text



%% The Appendices part is started with the command \appendix;
%% appendix sections are then done as normal sections
%% \appendix

%% \section{}
%% \label{}

%% If you have bibdatabase file and want bibtex to generate the
%% bibitems, please use
%%
\bibliographystyle{elsarticle-harv} 
\bibliography{bib}

%% else use the following coding to input the bibitems directly in the
%% TeX file.

%\begin{thebibliography}{00}

%% \bibitem[Author(year)]{label}
%% Text of bibliographic item

\bibitem[ ()]{}

%\end{thebibliography}
\end{document}

\endinput
%%
%% End of file `elsarticle-template-harv.tex'.
